Although the S+P transform is defined via local integer operations, its behavior is illuminating when viewed in the frequency domain. If we ignore truncation for the sake of analysis and treat the operations as linear, the mapping from the original sequence $c[n]$ to the highpass residual $d[n]$ can be represented as a linear time-invariant filter with a frequency response $F(\mathrm{e}^{j\omega})$.

The S stage alone already suppresses low frequencies in the highpass band, but it does so relatively mildly; the resulting highpass spectrum still exhibits considerable low-frequency energy due to the piecewise-constant nature of the Haar-like filter. By embedding prediction in the transform, S+P effectively implements an additional highpass filter: the predictor attempts to explain away low-frequency variations in $d_1[n]$ using neighboring lowpass differences and adjacent highpass values. When the predictor coefficients are chosen appropriately, $|F(\mathrm{e}^{j\omega})|$ exhibits significantly stronger attenuation at low frequencies than the basic S transform.

There is, of course, a trade-off. Stronger low-frequency attenuation generally comes at the cost of some amplification in higher frequencies. For smooth, low-noise images, this trade-off is supposedly advantageous: the additional whitening reduces the variance and first-order entropy of $d[n]$, improving compressibility. For highly textured or noisy images, it may be beneficial to temper the predictor so as not to over-amplify high-frequency content. Said and Pearlman demonstrate that, in practice, a small set of carefully tuned predictors can work well across a broad range of natural and medical images~[9]. 